The negative-cell premise says exactly that \(\omega\) is a sample from \(\mathsf W_j\) after exact specialization. Hence, by the definition of \(\mathsf W_j\),
\[
\mathsf C_{\mathcal G}(\theta^i;P_{\mathrm{obs}})
\quad(i=0,1),
\qquad
q_{a,y}(\theta^0)\neq q_{a,y}(\theta^1).
\tag{1}
\]
Applying \(\text{lem:negative-cell-model-pair}\) to (1) gives
\[
\mathcal M_i\in\mathfrak F(\mathcal G,P_{\mathrm{obs}}),
\qquad
P_{\mathcal M_i}(\mathcal O)=P_{\mathrm{obs}},
\qquad
Q_a^{\mathcal M_i}(y)=q_{a,y}(\theta^i)
\quad(i=0,1).
\tag{2}
\]
The same lemma proves that the decoded factor law is the product over the individual bidirected-edge and private factors and that every full-data cell is strictly positive. Substitution of the last equality in (2) into the inequality in (1) yields
\[
Q_a^{\mathcal M_0}(y)\neq Q_a^{\mathcal M_1}(y).
\tag{3}
\]

It remains only to check the compiler's return and termination claims. By the corrected \(\text{def:certificate-compiler}\), decoding is performed before the diagnostic scan. Every graphical predicate is decidable by inspection of the finite graph. Every trace equality or nonzero test is a sign test for a polynomial, after clearing the fixing denominators that are nonzero on the supported trace. The sample coordinates are represented by Thom encodings over the ordered field generated by \(\mathbf p\), and \(\operatorname{Exact}(P_{\mathrm{obs}})\) decides the remaining polynomial signs; hence each predicate test terminates. Since \(\mathcal T\) is finite, scanning it in deterministic prefix order performs only finitely many such tests. If one succeeds, the construction returns the candidate triple; if none succeeds, it returns the failure triple. In either branch the model components are exactly the decoded models in (2), so (2)--(3) prove every affirmative property. The corrected construction does not infer any graphical certificate property from the scan, which proves the statement without a hidden colored-failure invariant.